% Generated draft: requires mathematical and source review.
\documentclass[11pt]{article}
\usepackage[margin=25mm]{geometry}
\usepackage{fontspec}
\setmainfont{Latin Modern Roman}
\usepackage{amsmath,amssymb,mathtools,unicode-math}
\setmathfont{Latin Modern Math}
\usepackage{enumitem,needspace,xurl,hyperref}
\setlength{\parindent}{0pt}
\setlength{\parskip}{5pt}
\setlength{\emergencystretch}{3em}
\newcommand{\N}{\mathbb{N}}
\newcommand{\Z}{\mathbb{Z}}
\newcommand{\Q}{\mathbb{Q}}
\newcommand{\R}{\mathbb{R}}
\newcommand{\C}{\mathbb{C}}
\newcommand{\F}{\mathbb{F}}
\newcommand{\A}{\mathbb{A}}
\newcommand{\PP}{\mathbb{P}}
\newcommand{\E}{\mathbb{E}}
\newcommand{\eps}{\varepsilon}
\newcommand{\dd}{\,\mathrm{d}}
\begin{document}
\section*{UnsolvedMath ID: 701480}
\textbf{Source problem code:} catalogue entry
\subsection{Definitions and mathematical statement}
Let \(p\) be a prime. For each positive integer \(n\), define
\[
f_p(n)=\min\left\{m\in\Z_{\geq 0}:\ \exists\text{ a finite group }U,\ |U|=p^m,\ \forall\text{ finite groups }G,\ |G|=p^n\Longrightarrow\exists\text{ an injective group homomorphism }G\to U\right\}.
\]
Thus a single group \(U\) must contain an isomorphic copy of every group of order \(p^n\). The minimum exists because there are finitely many isomorphism classes of groups of a given finite order, and the direct product of representatives is a finite \(p\)-group containing every representative. If the source convention includes \(0\in\N\), the same definition gives \(f_p(0)=0\); this does not affect the asymptotic question.

The source asks whether \(f_p(n)\) grows faster than polynomially but slower than exponentially as \(n\to\infty\). Under the interpretation that \(p\) is fixed, that growth exceeds every polynomial, and that growth is subexponential, the question is whether, for every prime \(p\),
\[
\left(\forall k\in\R_{>0},\quad\lim_{n\to\infty}\frac{f_p(n)}{n^k}=+\infty\right)
\quad\text{and}\quad
\left(\forall a\in\R_{>1},\quad\lim_{n\to\infty}\frac{f_p(n)}{a^n}=0\right).
\]
The source does not explicitly specify these asymptotic quantifiers or any uniformity in \(p\); the displayed limits are an explicit interpretation requiring review, rather than additional hypotheses supplied by the source.
\subsection{Short English statement}
For each prime \(p\), find the smallest exponent of a finite \(p\)-group into which every group of order \(p^n\) embeds. Does this exponent eventually grow faster than every polynomial in \(n\), while remaining smaller than every exponential with base greater than \(1\), for fixed \(p\)? The source's informal growth terminology requires confirmation.
\subsection{Sources}
\begin{itemize}
\item Supplied local UnsolvedMath registry record, permanent ID \(701480\), exact source problem code \texttt{catalogue entry}, title \emph{Kourovka Notebook Problem 18.51}. The supplied registry labels the problem \texttt{open}; this label has not been independently verified.
\item \emph{The Kourovka Notebook (Unsolved Problems in Group Theory)}, original problem 18.51, 18th issue (2014), attributed to S. Kohl. The supplied background identifies the 21st issue edition (2026), arXiv identifier \texttt{1401.0300v48}. The collection description names E. I. Khukhro and V. D. Mazurov as editors.
\item The registry record was updated on 2026-10-16. This is supplied registry metadata, not a date of independent literature or source verification. No separate literature assessment or literature-source links were supplied.
\end{itemize}
\end{document}
